\documentclass[10pt,a4paper]{amsart}
\usepackage[margin=19mm]{geometry}
\usepackage{amsmath,amssymb,mathtools}
\usepackage[scr=rsfs,cal=euler]{mathalfa}
\usepackage{xfrac}
\usepackage[unicode,hidelinks,bookmarks=false]{hyperref}
\newcommand{\ie}{i.e.\ }
\newcommand{\cf}{cf.\ }
\newcommand{\resp}{respectively\ }
\newcommand{\Subsection}{\subsection}
\providecommand{\nobreakdash}{\nobreak\hbox{-}}
\setlength{\emergencystretch}{2em}
\multlinegap0pt
\usepackage{mathtools}
\usepackage{amssymb}
\usepackage{bm}
\usepackage{tikz}
\usepackage{siunitx}
\usepackage{algorithm}\usepackage{algpseudocode}
\newtheorem{lemma}{Lemma}
\title{Graph-Based Meshfree Multi-scale Coarse Space Approximation for Two-Level Schwarz Methods}
\author{Yucheng Liu, Tak Shing Au Yeung, Eric T. Chung, Simon See}
\date{Source: arXiv:2603.26149v1 (CC BY 4.0)}
\begin{document}
\maketitle

\noindent\emph{Source and licence.} Graph-Based Meshfree Multi-scale Coarse Space Approximation for Two-Level Schwarz Methods. Version arXiv:2603.26149v1 (CC BY 4.0), distributed by arXiv under the Creative Commons Attribution 4.0 International licence, http://creativecommons.org/licenses/by/4.0/, as recorded in the arXiv licence metadata for this exact version and shown on the abstract page https://arxiv.org/abs/2603.26149v1. Full licence notice: \texttt{licenses/arxiv-2603.26149-CC-BY-4.0.txt}.\par\smallskip

\noindent\textit{Editorial scope.} Counted unit: the article's lemma conddep (source0.tex lines 571--589). The article's complete original proof of it is delivered with it (source0.tex lines 590--620, one \texttt{proof} environment of 31 lines). No mathematics is written, repaired or re-derived: the statement and the proof are the article's own lines, in the article's own notation, and no retained line is substituted or re-flowed.  No further source-local context is needed: every cross-reference of every retained line resolves inside the two delivered spans.  Nothing was omitted from any retained span, so the package carries no omission record.  No retained line contains a citation, so the wrapper carries no bibliography and no bibliography entry.  No retained line contains a figure environment, an image-inclusion command or drawing code, so no image is delivered and the package declares no asset dependency.  Editorial formatting only, in addition: an AMS \texttt{amsart} preamble that restates the article's own declarations and macros used by the retained lines, namely the environments \texttt{lemma} on separate counters, together with the standard packages those lines need.  The retained environments print in this document as Lemma 1 in order of appearance, whereas the article prints the same items with the numbering of its own full document.  The complete native source of the article as distributed in the arXiv e-print for this exact version is bundled unchanged as \texttt{sources/197260/source0.tex}.
\begin{lemma}\label{conddep}
    Let \( B, \tilde{B} \in \mathbb{R}^{m \times m} \) be two symmetric positive semidefinite matrices.
    Let \( \ker(\tilde{B}) \), \( \Im(\tilde{B}) \) be the null space and image space of \( \tilde{B} \), respectively.
    Suppose subspaces $K$ and $I$ satisfy
\[
I \subseteq \Im(\tilde B), \quad (\ker(B)\cap \ker(\tilde B))^{\perp} \cap \ker(\tilde B) \subseteq K \subseteq \ker(\tilde B).
\]
Let $\Pi$ denote the orthogonal projector onto the subspace $S := K \oplus I$ and $\tau>0$ be arbitrary. Then the following two statements are equivalent.
\begin{itemize}
    \item For any $u \in \mathbb{R}^{m}$, we have
    \[
    (u - \Pi u)^\top B (u - \Pi u) \leq \tau u^\top \tilde B u.
    \]
    \item For any $v \in \Im(\tilde B)$, $\Pi_I$ is the $\tilde{B}$-orthogonal projector onto the subspace $I$, we have
    \[
    (v - \Pi_I v)^\top B (v - \Pi_I v) \leq \tau v^\top \tilde B v.
    \]
\end{itemize}
\end{lemma}

\begin{proof}
    Because \(\tilde{B}\) is symmetric positive semidefinite matrix, we have the orthogonal decomposition
    \[\Im(\tilde B) \oplus \ker(\tilde B) =   \mathbb{R}^{m}.\]
    \begin{itemize}
        \item [\((\Rightarrow)\)] Assume the first statement holds.
        Let \( v \in \Im(\tilde B) \subset \mathbb{R}^{m} \). Then \( v \) satisfies
        \[
        (v - \Pi v)^\top B (v - \Pi v) \leq \tau v^\top \tilde B v.
        \]
       By the definition of \( \Pi \) and \( S = K \oplus I \), the projection \( \Pi v \) admits a unique decomposition
       \[
       \Pi v = u_K + u_I, \quad \text{with } u_K \in K,\ u_I \in I.
       \]
       Since \( K \subseteq \ker(\tilde B) \), we have \( u_K \in \ker(\tilde B) \). But \( \Im(\tilde B) \perp \ker(\tilde B) \) and $v \in \Im(\tilde B)$, so \( u_K = 0 \). Hence, \( \Pi v = u_I = \Pi_I v\). We obtain
       \[
       (v - \Pi_I v)^\top B (v - \Pi_I v) \leq \tau v^\top \tilde B v.
       \]  
        \item[\((\Leftarrow)\)] Assume the second statement holds. Let \( u \in \mathbb{R}^{m} \). Write \( u = v + w \), where \( v \in \Im(\tilde B) \), \( w \in \ker(\tilde B) \). By assumption,
        \[
        (v - \Pi_I v)^\top B (v - \Pi_I v) \leq \tau v^\top \tilde B v.
        \]
        Due to \((\ker(B)\cap \ker(\tilde B))^{\perp} \cap \ker(\tilde B) = \ker^{\perp}(B)\cap\ker(\tilde B)\), we have \(\ker^{\perp}(B)\cap\ker(\tilde B) \subset K \subset \ker(\tilde B)\). By the definition of \( \Pi \) and $w \in \ker(\tilde B)$, we can get \(w - \Pi w \in \ker(\tilde B) \setminus K \subset \ker(\tilde B) \setminus (\ker^{\perp}(B)\cap\ker(\tilde B)) = \ker(\tilde B) \cap \ker(B)\). hence
        \[
        (u - \Pi u)^\top B (u - \Pi u) = (v - \Pi v)^\top B (v - \Pi v).
        \]
        Due to \(\Pi v = \Pi_I v\), 
        \[
        (u - \Pi u)^\top B (u - \Pi u) = (v - \Pi_I v)^\top B (v - \Pi_I v) \leq \tau u^\top \tilde B u.
        \]
    \end{itemize}
\end{proof}

\end{document}
